Put \(r=1/4\), and first consider a generic design in the construction whose early and late proxies have signals \(x,y\in(0,1)\), respectively. Conditional on \(U\), the signed treatment \(T\), the early proxy \(W_x\), and the late proxy \(W_y\) are mutually independent, with conditional means \(rU\), \(xU\), and \(yU\). Since \(U\) is Rademacher, marginalizing over \(U\) gives \(p_{g,a}=1/2\) in both sources.

We first verify regularity rather than taking membership in \(\mathfrak P_{\mathcal D}\) from the construction. Fix \(T=t\). The two-by-two matrix of \(T_d\), with rows indexed by the early proxy and columns indexed by the late proxy, has determinant
\[
 \det(T_d)=\frac{xy(1-r^2)}{1-r^2x^2}>0. \tag{1}
\]
Indeed, this determinant is the difference between the two conditional probabilities of \(W_y=1\), and Bayes' rule gives
\[
 \frac y2\left\{
 \frac{x+rt}{1+rtx}-\frac{-x+rt}{1-rtx}
 \right\}
 =\frac{xy(1-r^2)}{1-r^2x^2}.
\]
For a single proxy \(W_s\), the determinant of the conditional-expectation map from functions of \(U\) to functions of \(W_s\), conditional on \(T=t\), is similarly
\[
 \frac{s(1-r^2)}{1-r^2s^2}. \tag{2}
\]
It is nonzero for \(s=x\) and \(s=y\). Thus the early- and late-proxy completeness implications hold, while (1) proves designwise operator bijectivity.

The bridge formulas can be checked directly. Conditional independence gives
\[
 \mathbb E\!\left[\frac{W_y}{y}\mathrel{\big|}W_x,T,G=O\right]
 =\mathbb E[U\mid W_x,T,G=O]
 =\mathbb E[Y\mid W_x,T,G=O]. \tag{3}
\]
Hence \(h_d(l,a)=l/y\) solves the outcome-bridge equation. Moreover,
\[
 \frac{p(W_y=l\mid T=t,G=E)}{p(W_y=l\mid T=t,G=O)}
 =\frac1{1+rtyl}. \tag{4}
\]
Since
\[
 \mathbb E[W_x\mid W_y=l,T=t,G=O]
 =x\frac{rt+yl}{1+rtyl},
\]
the proposed selection bridge satisfies
\[
 \mathbb E\!\left[
 \frac{1-rtW_x/x}{1-r^2}\mathrel{\big|}W_y=l,T=t,G=O
 \right]
 =\frac1{1+rtyl}. \tag{5}
\]
This is exactly (4), so \(q_d(e,a)=\{1-r(2a-1)e/x\}/(1-r^2)\) solves the selection-bridge equation. Equation (1) also makes the outcome bridge unique.

Every primitive cell probability is positive because \(r,x,y\in(0,1)\). Therefore observational and experimental overlap and cross-source support hold. Consistency holds by the definitions of \(Y\) and \(Z\); treatment is independent of all potential variables given \(U\) in source \(O\), and it is independent of them unconditionally in source \(E\). The joint potential short-term vector and \((U,X)\) are independent of \(G\). Conditional on \(U\), the pair consisting of \(Y(a)=U\) and the late proxy is independent of the early proxy, so the sequential-outcome restriction holds. Finite support supplies square integrability and all needed moments. Together with (1)--(5) and the positive variances computed below, these facts verify all member conditions of \(\mathfrak P_{\mathcal D}\).

We record why the efficiency comparison is causal throughout the relevant constrained local model, rather than merely at the center \(P_\epsilon\). Let \(\widetilde P\) be any regular law in that constrained model, in which \(X\) and every \(S_{2d}\) are constant. For each \(d\) and \(a\), define the function of the latent variable
\[
 \delta_{d,a}(U)
 =\mathbb E_{\widetilde P}\!\left[
 Y(a)-h_d(S_{3d}(a),a)\mid U,G=O
 \right].
\]
Consistency, observational latent exchangeability, the sequential-outcome restriction, and the observed outcome-bridge equation imply
\[
 0
 =\mathbb E_{\widetilde P}\!\left[
 Y-h_d(S_{3d},a)\mid S_{1d},A=a,G=O
 \right]
 =\mathbb E_{\widetilde P}\!\left[
 \delta_{d,a}(U)\mid S_{1d},A=a,G=O
 \right]. \tag{6}
\]
Because \(X\) and \(S_{2d}\) are constant, the discrepancy in (6) is a function of \(U\) alone. Early-proxy completeness therefore gives \(\delta_{d,a}(U)=0\) almost surely. Experimental randomization and short-term consistency, followed by external validity and the preceding equality, now yield
\[
 \begin{aligned}
 \mu_{d,a}
 &=\mathbb E_{\widetilde P}[h_d(S_{3d}(a),a)\mid G=E]\\
 &=\mathbb E_{\widetilde P}[h_d(S_{3d}(a),a)\mid G=O]\\
 &=\mathbb E_{\widetilde P}[Y(a)\mid G=O].
 \end{aligned} \tag{7}
\]
Thus \(\mu_{d,1}-\mu_{d,0}=\tau\) for both designs throughout this local identified submodel. In particular, at \(P_\epsilon\), \(Y(1)=Y(0)=U\), so \(\tau=0\), and symmetry also gives \(\mu_{d,a}=\mathbb E[W_y/y\mid A=a,G=E]=0\). By the fixed-system efficiency lemma, bridge uniqueness and (1) make the source gradients canonical for this causal target on the constrained local submodel. Hence the profiled values compared below are causal efficiency bounds there, not merely variances of unrelated bridge functionals.

We next calculate the two objectives at \(P_\epsilon\). The best predictor of \(U\) from \((W_x,W_y)=(e,l)\) is
\[
 \mathbb E[U\mid W_x=e,W_y=l]
 =\frac{xe+yl}{1+xyel}.
\]
Averaging its conditional variance over the four proxy cells gives
\[
 R(x,y)=\frac{(1-x^2)(1-y^2)}{1-x^2y^2}. \tag{8}
\]
Consequently,
\[
 R_{d_{\mathrm p}}(P_\epsilon)
 =\frac{7(1-\epsilon^2)}{16-9\epsilon^2},
 \qquad
 R_{d_{\mathrm e}}(P_\epsilon)=\frac35. \tag{9}
\]
The paired-source costs are \(6\) and \(5\), respectively. Therefore
\[
 B\Psi_{\mathrm{CLW},d_{\mathrm p}}(P_\epsilon;B)
 =\frac{42(1-\epsilon^2)}{16-9\epsilon^2}
 <3
 =B\Psi_{\mathrm{CLW},d_{\mathrm e}}(P_\epsilon;B), \tag{10}
\]
where the strict inequality is equivalent to \(42(1-\epsilon^2)<48-27\epsilon^2\). Thus \(d_{\mathrm p}\) is the unique predictive optimizer despite its larger cost.

For the efficiency criterion, the experimental score reduces to \(2TW_y/y\), and hence
\[
 V_E(x,y)=\frac4{y^2}. \tag{11}
\]
The observational score is
\[
 2Tq_d(W_x,A)\left(U-\frac{W_y}{y}\right).
\]
Conditional on \(U\), the squared residual has expectation \((1-y^2)/y^2\) and is independent of \((T,W_x)\). In addition, \(\mathbb E[TW_x\mid G=O]=rx\), so
\[
 \mathbb E[q_d(W_x,A)^2\mid G=O]
 =\frac{1-2r^2+r^2/x^2}{(1-r^2)^2}.
\]
It follows that
\[
 V_O(x,y)=
 \frac4{(1-r^2)^2}\frac{1-y^2}{y^2}
 \left(1-2r^2+\frac{r^2}{x^2}\right). \tag{12}
\]
Equations (11)--(12) are strictly positive and finite for every fixed \(x,y\in(0,1)\), completing the verification of nondegeneracy.

For \(d_{\mathrm p}\), substituting \((x,y)=(\epsilon,3/4)\) into (11)--(12) gives
\[
 V_{E,d_{\mathrm p}}=\frac{64}{9},
 \qquad
 V_{O,d_{\mathrm p}}
 =\frac{448}{2025\epsilon^2}+\frac{6272}{2025}. \tag{13}
\]
For \(d_{\mathrm e}\), substituting \((x,y)=(1/2,1/2)\) gives
\[
 V_{E,d_{\mathrm e}}=16,
 \qquad
 V_{O,d_{\mathrm e}}=\frac{384}{25}. \tag{14}
\]
The budget-profile theorem and the costs \(c_{E,d_{\mathrm e}}=2\), \(c_{O,d_{\mathrm e}}=3\) now imply
\[
 \mathcal V_{d_{\mathrm e}}^*(P_\epsilon)
 =\left(\sqrt{2\cdot16}+\sqrt{3\cdot\frac{384}{25}}\right)^2
 =\left(4\sqrt2+\frac{24\sqrt2}{5}\right)^2
 =\frac{3872}{25}. \tag{15}
\]
On the other hand, for \(0<\epsilon<10^{-3}\),
\[
 \mathcal V_{d_{\mathrm p}}^*(P_\epsilon)
 \ge c_{O,d_{\mathrm p}}V_{O,d_{\mathrm p}}
 >\frac{1568}{2025\cdot10^{-6}}
 >\frac{3872}{25}. \tag{16}
\]
Since \(\mathcal D\) contains only the two designs, (16) proves that \(d_{\mathrm e}\) is the unique design-value optimizer. Expanding the squared profile using (13) yields
\[
 \mathcal V_{d_{\mathrm p}}^*(P_\epsilon)
 =\frac{1568}{2025\epsilon^2}+O(\epsilon^{-1}). \tag{17}
\]
Dividing (17) by (15), and using \(1568=32\cdot49\) and \(3872=32\cdot121\), proves
\[
 \rho_{\mathrm{pred}}(P_\epsilon)
 =\frac{49}{9801\epsilon^2}+O(\epsilon^{-1}). \tag{18}
\]

It remains to prove stability. Fix \(\epsilon\) and parameterize the constrained fixed-support model by its finitely many free latent-cell probabilities. In a sufficiently small neighborhood of \(P_\epsilon\), all conditioning denominators remain positive. Conditional-probability matrices and their determinants are continuous functions of the cell probabilities. Wherever those determinants are nonzero, the bridge coefficients obtained by finite-dimensional inversion are rational, hence continuous. The prediction risks and source-score second moments are finite sums of continuous functions of the cell probabilities and bridge coefficients. The profiled values are continuous because they are compositions of these positive quantities with square roots.

At \(P_\epsilon\), (1)--(2), (10), and (16) give strictly positive regularity, predictive-selection, and design-value-selection margins. Continuity therefore furnishes a radius \(r_\epsilon>0\) on which both determinants remain nonzero and the same two strict optimizer assignments persist. The denominator of \(\rho_{\mathrm{pred}}\) is positive by (14)--(15), so the ratio is continuous as well. Shrinking \(r_\epsilon\) if necessary gives
\[
 \rho_{\mathrm{pred}}(P)>\frac12\rho_{\mathrm{pred}}(P_\epsilon)
\]
throughout the stated relative ball. This proves every assertion of the lemma.